\section{Problem Statement}

This preprint addresses exactly one scope-locked question:

\begin{quote}
\emph{Q5: What are the formal limits of KPI-based enterprise analysis in
continuum-modeled systems \(K_{EA}\)?}
\end{quote}

\paragraph{Hard constraints.}
All statements are derived strictly within the Ontology of Continua (OC) and the
executable theory specification \texttt{ETS.K\_EA.v1.1}.
No new primitives are introduced, and no reinterpretation or extension of OC or
ETS is permitted.
The scope is purely diagnostic: the results are negative in nature and establish
existence and impossibility claims concerning what KPI-based analysis cannot
identify.

\paragraph{Explicit non-claims.}
This work makes no managerial, governance, optimization, or control claims.
It does not recommend KPIs, prescribe dashboards, propose decision rules, or
advance performance narratives.
No probabilistic or statistical enrichment is introduced beyond what is already
contained in OC/ETS.
References to aggregation, sampling, or post-processing are included solely to
emphasize that the derived limits are structural and do not disappear under such
operations.

\paragraph{Core distinction.}
The central distinction formalized in this work is between
\emph{observability} and \emph{identifiability}.
Observability concerns whether readouts are defined and admissible on a given
domain, whereas identifiability concerns whether those readouts uniquely
determine the underlying system state.
This distinction is classical in systems theory and system identification
\cite{kalman1960,ljung1999,walter1997}.
While OC/ETS explicitly provides observability structures at the level of the
enterprise continuum, the question addressed here is whether KPI-based readouts
can be identifying in the formal sense.
The main result is that, in continuum-modeled enterprise systems,
KPI-based observation is generically non-identifying.

\paragraph{Formal relevance.}
Under OC/ETS, enterprises are modeled as continua with state space
\(\Omega(K_{EA})\), structured by axes, flows, thresholds, cycles, and a
continuum measure.
KPI readouts are, by construction, finite-dimensional projections of this state
space.
It is a standard result in identification theory that such projections generally
induce equivalence classes of indistinguishable states unless injectivity is
explicitly enforced \cite{ljung1999,walter1982}.
This preprint formalizes that reduction in the OC/ETS setting and shows that it
induces unavoidable non-identifiability: multiple structurally distinct
enterprise states can share identical KPI values.
